% Step 9 historical addendum, substantially corrected in Step 10.
% The claims below that Taub--NUT--AdS refutes adjustable-buffer cover
% regularity, has an all-radius bounded curvature maximal function, or forces
% an exclusively global-topology route are withdrawn.  The exact correction is
% workspace/notes/step10_multiscale_cover_and_diffuse_pontryagin.tex.

\section{Superseded assessment of Taub--NUT--AdS}

\textbf{Step--10 correction.}  The computation below only controls balls of
radius at most $1/3$ around a point at distance $1/2$ from the NUT.  Balls
reaching the NUT have divergent normalized curvature energy as $\nu\to0$.
Moreover, choosing a smaller buffered domain before taking its universal
cover unwraps the Hopf circle at a fixed scale.  The unwrapped annular limit
is real hyperbolic and has calibrated scalar--Weyl deficit $16$, not $0$.
Thus this family does not refute the simultaneous multiscale
curvature-plus-calibration accessibility question.  See the Step--10 note for
exact Weyl formulas, corrected volume factors, and homothetic scaling.

\textbf{The claimed regularity}---that controlled maximal $\fint|\operatorname{Rm}|^2 \le K$
at every scale on the local universal cover, together with $\operatorname{Ric}\ge-3g$,
forces noncollapse---is \textbf{false}.

\subsection{The countermodel}
The Taub--NUT--Anti-de~Sitter metric with NUT parameter $\nu>0$
and cosmological length $\ell=\sqrt3$ is
\[
g_\nu = \frac{dr^2}{F_\nu(r)}
+ (r^2-\nu^2)(d\theta^2+\sin^2\theta\,d\phi^2)
+ 4\nu^2 F_\nu(r)(d\psi+\cos\theta\,d\phi)^2,
\]
where
\[
F_\nu(r) = \frac{r-\nu}{r+\nu}
\left(1+\frac{(r-\nu)(r+3\nu)}{\ell^2}\right),
\qquad r\ge\nu.
\]
This metric satisfies \textbf{exactly}
\[
\operatorname{Ric}_{g_\nu} = -g_\nu.
\]
At $r=\nu$ the metric has a smooth ``NUT'' (bolt of dimension zero):
after the coordinate change $r=\nu+\rho^2/(8\nu)$,
the metric near $\rho=0$ is the flat $\mathbb R^4$ metric.

\subsection{Why it is a countermodel}
Choose a basepoint $p_\nu$ at distance $1/2$ from the NUT.
For all $q\in B_{1/3}(p_\nu)$, the distance to the NUT lies in $[1/6,5/6]$.
In this fixed annular band:
\begin{enumerate}
\item $\sup|\operatorname{Rm}|^2 \le K_0$ (curvature bounded uniformly in $\nu$).
\item Hence $\fint_{B_r(p_\nu)}|\operatorname{Rm}|^2 \le K_0$ for all $r\le 1/3$.
\item The volume form scales as $\nu\cdot\sin\theta\,(r^2-\nu^2)\,dr\,d\theta\,d\phi\,d\psi$,
so $\operatorname{Vol}(B_s(p_\nu)) = O(\nu) \to 0$ for fixed $s$.
\end{enumerate}

\subsection{Why the local universal cover is still collapsed}
The $\psi$-circle at $p_\nu$ has length $O(\nu)$ and generates collapse.
However, this circle can be \emph{contracted to a point}
by sweeping radially to the NUT at $r=\nu$, which lies
at distance $1/2 < 1$ from $p_\nu$.
Therefore the $\psi$-circle is \textbf{nullhomotopic in $B_1(p_\nu)$}.

The local universal cover of $B_1(p_\nu)$ does not unwrap
nullhomotopic loops. So the cover is isometric to the original
ball near the annular band, and
\[
\operatorname{Vol}(B_s(\tilde p_\nu)) = O(\nu) \to 0.
\]

\subsection{The precise obstruction}
The failure of the Naber--Zhang package is at the
\textbf{maximal fibered nilpotency rank} condition.
The GH model (dimension 3), its noncollapsing radius, and the
curvature bounds are all satisfied. But the short fiber
is contractible in the buffered domain, so its fundamental-group
image is trivial. The fibered rank is 0, not the required $4-3=1$.

\textbf{Maximal rank is an independent geometric hypothesis, not a
consequence of curvature bounds, Ricci pinching, or
controlled maximal curvature function.}

\subsection{What the countermodel does NOT settle}
\begin{itemize}
\item The Taub--NUT--AdS space is \emph{noncompact}.
On a closed manifold, global topology constrains which fibers
can contract. The Pontryagin class $p_1(\Lambda^2_+)$ integrates to
$2\chi+3\tau \ne 0$ on the critical Ishida--LeBrun manifold,
and the TTCS generator shows this must localize on regions
where the SD bundle is not pulled back from the base.
\item The total four-index Riemann energy is infinite because the Einstein
background has infinite volume, but its full-Weyl core energy is finite and
converges to $32\pi^2$ as $\nu\to0$.  This distinction matters for any
energy comparison.
\item The model has no Seiberg--Witten monopole classes.
The Ishida--LeBrun manifolds have specific monopole structure
that constrains the near-saturation energy profile.
\end{itemize}

\section{Revised research direction}

The general regularity approach (maximal curvature function $\to$ noncollapse)
is \textbf{not viable}. The research must now use the specific structure
of the closed manifold problem:

\subsection{What the form analysis establishes}
The TTCS generator analysis provides rank-dependent results:
\begin{itemize}
\item \textbf{Flat rank $\ge 2$:}
Flat regular torus bundles are excluded.
If $H$ is a definite pair (two independent near-harmonic forms with locally
nondegenerate Gram matrix), then closedness and constant amplitude
force $\nabla H \to 0$, contradicting the Weitzenb\"ock gradient energy
$\int|\nabla H|^2 = (1+o(1))\int|H|^2$.
\item \textbf{Flat rank 3:}
Three independent forms with locally nondegenerate Gram matrix
give a full $\Lambda^2_+$ trivialization with $\nabla \to 0$,
same contradiction. But this requires \emph{local} nondegeneracy
of the Gram matrix, not just global independence.
\item \textbf{Rank 1:}
A single form reduces to a weighted harmonic 1-form on the
3D base, which can have $O(1)$ horizontal rotation. Not excluded.
\item \textbf{Nil-twisted models:}
Heisenberg-type structure can produce constant-length, closed,
self-dual forms with $O(1)$ covariant derivative.
Not excluded by form properties alone.
\end{itemize}

\subsection{The Pontryagin localization mechanism}
The generator identifies $p_1(\Lambda^2_+) = 2\chi + 3\tau$ as a
localization tool. On a regular fibered collapse where the
self-dual forms are basic (pulled back from the base),
$\Omega_E \wedge \Omega_E = 0$ (the base has dimension $\le 3$).
So the nonzero $\int p_1 = 2A/3$ must concentrate on:
\begin{itemize}
\item Non-basic vertical connection curvature
\item Nil/bundle twist regions
\item Singular fibers or curvature concentration
\item Genuinely 4D (noncollapsed) regions
\end{itemize}
This gives a \emph{charge localization} mechanism:
a positive amount of topological energy must sit
in the ``bad'' regions that escape the regularity framework.

\subsection{New research target}
The problem is now:
\[
\boxed{
\text{Can a closed critical Ishida--LeBrun 4-manifold support a }
\delta_i\to0\text{ sequence}
}
\]
\[
\boxed{
\text{whose collapsed regions are entirely composed of
NUT-type contractible fibers,}
}
\]
\[
\boxed{
\text{nil/solvable cusps, and curvature-concentrating
singular strata?}
}
\]

The Pontryagin charge $2A/3$, the monopole near-saturation,
the Weitzenb\"ock gradient energy, and the $\int|\operatorname{Rm}|^2$
budget must together account for the energy distribution
across these regions. The question is whether the energy
budget is sufficient.

\section{Files}
\begin{itemize}
\item TTCS verifier (countermodel): \texttt{trace/ttcs/results/step9\_ttcs01\_verifier1\_integration.md}
\item TTCS generator (form analysis): \texttt{trace/ttcs/results/step9\_ttcs01\_generator1\_integration.md}
\item Verification script: \texttt{workspace/scripts/verify\_weitzenbock\_clean.py}
\end{itemize}
